% Clase 7 --- Por qué Gauss NO alcanza para sigma_interna = 0 (§2.2).
% (a) Configuración que Gauss no prohíbe: carga neta interna cero, pero con
%     zonas + y - que se cancelan y líneas de campo cruzando la cavidad.
% (b) Lo que la descarta es la ESTABILIDAD: esas cargas sentirían el campo y se
%     moverían, así que no puede ser un equilibrio.
\begin{center}
\begin{tikzpicture}[scale=1]

  % =============== (a) lo que Gauss permite ===============
  \begin{scope}
    \draw[figline, fill=figgray!14] (0,0) circle (2.0);
    \draw[figline, fill=white] (0,0) circle (1.05);

    % densidades opuestas en la cara interna
    \foreach \a in {105,130,...,255} {
      \node[figlbl,figred,scale=0.7] at (\a:1.05) {$+$};
    }
    \foreach \a in {-75,-50,...,75} {
      \node[figlbl,figblue,scale=0.7] at (\a:1.05) {$-$};
    }

    % líneas cruzando la cavidad
    \foreach \y in {-0.45,0,0.45} {
      \draw[figvecres] (-0.85,\y) -- (0.85,\y);
    }
    \node[figlblaux, anchor=north] at (0,-2.15) {$\vec E \neq 0$ en la cavidad};

    \node[figlbl, align=center] at (0,-3.15)
      {$Q_{\text{int}} = 0$: \ Gauss \textbf{no} lo prohíbe};
    \node[figlbl] at (0,-4.1) {(a) sólo con Gauss};
  \end{scope}

  % =============== (b) lo que descarta la estabilidad ===============
  \begin{scope}[xshift=7.2cm]
    \draw[figline, fill=figgray!14] (0,0) circle (2.0);
    \draw[figline, fill=white] (0,0) circle (1.05);
    \node[figlblaux] at (0,0.25) {$\vec E = 0$};
    \node[figlblaux] at (0,-0.3) {$\sigma_{\text{int}} = 0$};

    % las cargas se moverían: flechas de migración
    \draw[figvecres] (120:1.15) arc (120:65:1.15);
    \draw[figvecres] (240:1.15) arc (240:295:1.15);
    \node[figlblaux, anchor=south east] at (-2.15,1.65) {se moverían};

    \node[figlbl, align=center] at (0,-3.15)
      {en equilibrio \textbf{no} puede haber campo\\ donde hay portadores libres};
    \node[figlbl] at (0,-4.1) {(b) con Gauss \emph{y} estabilidad};
  \end{scope}

\end{tikzpicture}\\[0.5em]
{\small\itshape Gauss sólo dice que la carga \textbf{neta} de la cara interna es
cero, y eso lo cumple también el panel (a), con zonas $+$ y $-$ que se
compensan. Hace falta un segundo ingrediente: si hubiera esas densidades habría
campo en la cavidad, las cargas de la superficie se moverían y no sería un
equilibrio. \textbf{Gauss no alcanza; hay que agregar la estabilidad.}}
\end{center}
